\chapter{Gaussian closure and the LMMSE floor}
\label{ch:gaussian-closure-lmmse}

The paper's Proposition 1 is short: run the sum-product recursion of
\cref{ch:factor-graphs-bp}, but after every transition integral replace the message by
the Gaussian with the same mean and variance; then, whatever the innovation law, the
posterior means you obtain are those of the \emph{Gaussian} chain with the same
correlation and innovation variance, and those are the linear minimum-mean-square
estimator $e^{-t}\Sig_0\Sig_t^{-1}\vx$. The proposition turns an approximation that
practitioners make for convenience (Gaussian message passing, assumed-density
filtering, the Kalman smoother applied to a non-Gaussian model) into an object with a
name and a closed form, and it gives the paper an apparatus check that is used in
every experiment: a Gaussian-innovation kernel fitted by maximum likelihood must land
on that closed form. This chapter defines the projection, derives the linear
estimator from first principles, shows that the Gaussian chain's forward--backward is
the Kalman smoother, proves the proposition in full, derives the two consequences the
paper draws from it, and explains what the gap between the true score and the
Gaussian one measures.

\section{The Gaussian moment projection}

\begin{definition}[Gaussian moment projection]
\label{def:gc-proj}
For a non-negative integrable function $g$ on $\Reals$ with finite second moment,
let $\bar g=g/\!\int g$ be its normalisation, $m=\int a\,\bar g(a)\dd a$ its mean and
$v=\int(a-m)^2\bar g(a)\dd a$ its variance. The \term{Gaussian moment projection} of $g$
is $\Pgauss g=\normal{m}{v}$, the Gaussian density with the same first two moments as
$\bar g$.
\end{definition}

The operator is defined on shapes: the normalisation of $g$ is discarded, which is
harmless because messages are defined up to a constant. Two facts make $\Pgauss$ the
natural projection rather than one choice among many. First, among all Gaussians,
$\Pgauss g$ is the one closest to $\bar g$ in Kullback--Leibler divergence
$\KL(\bar g\,\|\,\cdot)$: writing $\normal{m'}{v'}$ for a candidate,
$-\int\bar g\log\normal{m'}{v'}=\tfrac12\log(2\pi v')+\bigl(v+(m-m')^2\bigr)/2v'$ is
minimised at $m'=m$ and $v'=v$, since the term is a function of the first two moments of
$\bar g$ only. Second, and this is the property the proof of Proposition 1 uses, a
Gaussian is determined by its first two moments, so $\Pgauss$ forgets everything about
$g$ except $(m,v)$.

\term{Moment-projected sum-product} is the recursion of \cref{ch:factor-graphs-bp} with
$\Pgauss$ applied to each message after every transition integral: in the notation of
that chapter,
\begin{equation}
\begin{aligned}
  \lmsg^{\mathrm p}_i&=\Pgauss\!\Bigl[\int\kernel(\cdot\mid a')\like(x_{i-1}\mid a')\lmsg^{\mathrm p}_{i-1}(a')\dd a'\Bigr],\\
  \rmsg^{\mathrm p}_i&=\Pgauss\!\Bigl[\int\kernel(a''\mid\cdot)\like(x_{i+1}\mid a'')\rmsg^{\mathrm p}_{i+1}(a'')\dd a''\Bigr],
\end{aligned}
  \label{eq:gc-projected}
\end{equation}
with the same starting values as before. In the literature this is Gaussian
assumed-density filtering, or expectation propagation with a single Gaussian per
message and one forward and one backward pass. The paper's phrase is ``single-Gaussian
message closure''. Note where the projection is applied: after the transition integral,
and not after the multiplication by the unary factor, which would be a different
algorithm. For the Gaussian likelihood of the channel the difference is immaterial,
because a Gaussian message times a Gaussian likelihood is already Gaussian; the paper
fixes the order once so that the statement is unambiguous.

\section{The linear minimum-mean-square estimator}

\subsection{Derivation}

Let $(\va,\vx)$ be a pair of random vectors with finite second moments, means zero,
and second-moment matrices $\Sig_{aa}=\E[\va\va^{\!\top}]$, $\Sig_{ax}=\E[\va\vx^{\!\top}]$,
$\Sig_{xx}=\E[\vx\vx^{\!\top}]$, the last invertible. A \term{linear estimator} of $\va$
from $\vx$ is $\hat\va=G\vx$ for a fixed matrix $G$, and its \term{mean-square error} is
$R(G)=\E\lVert\va-G\vx\rVert^2=\tr\Sig_{aa}-2\tr(G\Sig_{ax}^{\!\top})+\tr(G\Sig_{xx}G^{\!\top})$,
by expanding the square and using linearity of expectation. The function $R$ is a
quadratic in the entries of $G$ with positive semidefinite Hessian (it is
$\tr(G\Sig_{xx}G^{\!\top})$ plus lower-order terms, and $\Sig_{xx}\succ0$), so its
stationary point is its minimum. Differentiating, $\partial R/\partial G=-2\Sig_{ax}+2G\Sig_{xx}$,
which vanishes at
\begin{equation}
  G^\ast=\Sig_{ax}\Sig_{xx}^{-1},\qquad
  \lmmse=\Sig_{ax}\Sig_{xx}^{-1}\vx,\qquad
  R(G^\ast)=\tr\bigl(\Sig_{aa}-\Sig_{ax}\Sig_{xx}^{-1}\Sig_{ax}^{\!\top}\bigr).
  \label{eq:gc-lmmse}
\end{equation}
This is the \term{linear minimum-mean-square estimator} (LMMSE), also called the Wiener
or normal-equation estimator; the matrix $G^\ast$ solves the \term{normal equations}
$G\Sig_{xx}=\Sig_{ax}$. The error $\va-G^\ast\vx$ is uncorrelated with $\vx$ (the
\term{orthogonality principle}: $\E[(\va-G^\ast\vx)\vx^{\!\top}]=\Sig_{ax}-G^\ast\Sig_{xx}=0$),
which is another way of saying that no linear correction can lower the error.

Apply this to the channel $\vx=e^{-t}\va+\sqrt{\Dt}\vz$ with $\va\sim\Pz$ of covariance
$\Sig_0$ and $\vz\sim\normal{0}{\Id}$ independent of $\va$. Then
$\Sig_{ax}=\E[\va(e^{-t}\va+\sqrt{\Dt}\vz)^{\!\top}]=e^{-t}\Sig_0$ and
$\Sig_{xx}=e^{-2t}\Sig_0+\Dt\Id=\Sig_t$, so
\begin{equation}
  \lmmse=e^{-t}\Sig_0\Sig_t^{-1}\vx,\qquad \Sig_t=e^{-2t}\Sig_0+\Dt\Id.
  \label{eq:gc-lmmse-channel}
\end{equation}
The derivation used only the first two moments of $\Pz$: every law with covariance
$\Sig_0$ has the same LMMSE, which is why the paper calls the Gaussian with covariance
$\Sig_0$ the \term{covariance twin} of $\Pz$ and \cref{eq:gc-lmmse-channel} ``the LMMSE
of the covariance twin''.

\subsection{When it is the MMSE, and when it is only the best linear map}

The \term{minimum-mean-square estimator} without the word ``linear'' is the conditional
mean $\E[\va\mid\vx]$, the minimiser of $\E\lVert\va-f(\vx)\rVert^2$ over all measurable
$f$ (the standard argument: for any $f$, $\E\lVert\va-f\rVert^2=\E\lVert\va-\E[\va\mid\vx]\rVert^2+\E\lVert\E[\va\mid\vx]-f\rVert^2$,
the cross term vanishing by the tower property). If $(\va,\vx)$ is jointly Gaussian,
the conditional mean is linear, $\E[\va\mid\vx]=\Sig_{ax}\Sig_{xx}^{-1}\vx$, so the LMMSE
\emph{is} the MMSE. This is the case of the Gaussian chain: $\va$ Gaussian and $\vz$
Gaussian independent make $(\va,\vx)$ jointly Gaussian. For any other $\Pz$ the
conditional mean is a nonlinear function of $\vx$, the LMMSE is only the best among
linear maps, and the difference
$\E\lVert\E[\va\mid\vx]-\lmmse\rVert^2=\E\lVert\va-\lmmse\rVert^2-\E\lVert\va-\E[\va\mid\vx]\rVert^2$
is the price of linearity. Through Tweedie's identity, the score of the covariance
twin is
\begin{equation}
  \Sgauss(\vx,t)=-\frac{\vx-e^{-t}\lmmse}{\Dt}
  =-\frac{1}{\Dt}\bigl(\Id-e^{-2t}\Sig_0\Sig_t^{-1}\bigr)\vx=-\Sig_t^{-1}\vx,
  \label{eq:gc-score-twin}
\end{equation}
where the last step uses $\Id-e^{-2t}\Sig_0\Sig_t^{-1}=(\Sig_t-e^{-2t}\Sig_0)\Sig_t^{-1}=\Dt\Sig_t^{-1}$.
So $\Sgauss=-\Sig_t^{-1}\vx$ is at once the score of $\normal{0}{\Sig_t}$, which is the
noised law of the Gaussian twin, and the Tweedie image of the LMMSE. The paper's
Section 2 states both forms.

\section{The Gaussian chain and the Kalman smoother}
\label{sec:gc-kalman}

For the Gaussian chain, $\mu=\normal{0}{1}$, $\kernel(a'\mid a)=\normal{\corr a}{\ivar}$,
$\like(x\mid a)\propto\normal{x}{\Dt}$ evaluated at $e^{-t}a$, every message of
\cref{ch:factor-graphs-bp} is a Gaussian in $a$ and is carried by two numbers. It is
convenient to use the \term{information form}: a Gaussian shape $\propto\exp(-\lambda a^2/2+ha)$
has precision $\lambda$, information $h$, mean $h/\lambda$ and variance $1/\lambda$;
products of shapes add $(\lambda,h)$; a flat message is $\lambda=0$, $h=0$. The unary
factor, as a function of $a$, has $\lambda_{\like}=e^{-2t}/\Dt$ and
$h_{\like}=e^{-t}x_i/\Dt$.

\paragraph{Forward step.} Let $\lmsg_{i-1}=\normal{m}{v}$. Multiplying by $\like(x_{i-1}\mid\cdot)$
gives the \term{tilted} Gaussian with $\lambda'=1/v+\lambda_{\like}$, $h'=m/v+h_{\like}$,
mean $m'=h'/\lambda'$ and variance $v'=1/\lambda'$. The transition integral is
$\int\normal{\corr a'}{\ivar}(a)\,\normal{m'}{v'}(a')\dd a'$, the density of
$\corr a'+\varepsilon$ with $a'\sim\normal{m'}{v'}$ and $\varepsilon\sim\normal{0}{\ivar}$
independent, which is $\normal{\corr m'}{\corr^2v'+\ivar}$. Hence
\begin{equation}
  (m,v)\ \xrightarrow{\ \text{tilt}\ }\ (m',v')\ \xrightarrow{\ \text{transition}\ }\
  \bigl(\corr m',\ \corr^2v'+\ivar\bigr).
  \label{eq:gc-forward-map}
\end{equation}
This is the Kalman filter: the tilt is the measurement update and the transition is
the prediction step.

\paragraph{Backward step.} Let $\rmsg_{i+1}=\normal{m}{v}$ and tilt it by
$\like(x_{i+1}\mid\cdot)$ to $\normal{m'}{v'}$. The backward integral is
$\int\normal{\corr a}{\ivar}(a'')\,\normal{m'}{v'}(a'')\dd a''$, a function of $a$: it is
$h(\corr a)$ with $h(u)=\int\normal{u}{\ivar}(a'')\normal{m'}{v'}(a'')\dd a''$ the density
of $a''-\varepsilon$ at $u$, which is $\normal{m'}{v'+\ivar}(u)$. As a function of $a$,
$\normal{m'}{v'+\ivar}(\corr a)\propto\exp[-(\corr a-m')^2/2(v'+\ivar)]$ is a Gaussian
shape in $a$ with mean $m'/\corr$ and variance $(v'+\ivar)/\corr^2$, provided $\corr\ne0$:
\begin{equation}
  (m,v)\ \xrightarrow{\ \text{tilt}\ }\ (m',v')\ \xrightarrow{\ \text{transition}^{-1}\ }\
  \bigl(m'/\corr,\ (v'+\ivar)/\corr^2\bigr).
  \label{eq:gc-backward-map}
\end{equation}
At $\corr=0$ the backward message is constant in $a$, has no normalised moments, and
the map is undefined; the recursion is still fine (a flat message is $\lambda=0$), but
the moment description is not. This is the origin of the condition $\corr\ne0$ in the
proposition. The terminal message $\rmsg_\nsites\equiv1$ is flat and is never projected;
the backward induction starts from the tilt by $\like(x_\nsites\mid\cdot)$.

\paragraph{Beliefs.} At site $i$ the belief has precision
$\lambda_i^{\lmsg}+\lambda_{\like}+\lambda_i^{\rmsg}$ and information
$h_i^{\lmsg}+h_{\like}+h_i^{\rmsg}$; its mean is the posterior mean and its variance the
posterior variance. Combining a forward filter with a backward pass in this way is the
Rauch--Tung--Striebel smoother written in information form, and the code is
\code{gaussian\_chain\_bp} in \code{src/markov\_score/bp\_gaussian.py}, which is tested
against the closed form $e^{-t}\Sig_0\Sig_t^{-1}\vx$ to machine precision
(\code{tests/test\_gaussian\_bp\_equivalence.py}).

\section{Proposition 1 and its proof}

\begin{proposition}[Gaussian closure computes the LMMSE of the covariance twin]
\label{prop:gc-lmmse}
Let the chain have $a_1\sim\normal{0}{1}$, $a_i=\corr a_{i-1}+\varepsilon_i$ with
$\corr\ne0$ and innovations $\varepsilon_i$ i.i.d.\ from any law $\innov$ of mean zero
and variance $\ivar$. Then moment-projected sum-product \cref{eq:gc-projected} on the
chain posterior returns, at every site, the belief of the covariance-matched Gaussian
chain, and hence the posterior means $e^{-t}\Sig_0\Sig_t^{-1}\vx$, the same for every
such $\innov$.
\end{proposition}

The proof needs one lemma about convolutions.

\begin{lemma}[First two moments of a convolution]
\label{lem:gc-conv}
Let $U$ and $V$ be independent with finite variances, means $m_U,m_V$ and variances
$v_U,v_V$. Then $\corr U+V$ has mean $\corr m_U+m_V$ and variance $\corr^2v_U+v_V$, and
$U-V$ has mean $m_U-m_V$ and variance $v_U+v_V$.
\end{lemma}
\begin{proof}
Linearity of expectation gives the means. For the variances,
$\Var(\corr U+V)=\corr^2\Var U+\Var V+2\corr\Cov(U,V)$ and the covariance vanishes by
independence; likewise for $U-V$. No other property of the laws of $U$ and $V$ enters.
\end{proof}

\begin{proof}[Proof of \cref{prop:gc-lmmse}]
Write $\lmsg^{\mathrm p}_i$ and $\rmsg^{\mathrm p}_i$ for the projected messages and
$\lmsg^{\mathrm G}_i$, $\rmsg^{\mathrm G}_i$ for the exact messages of the Gaussian chain
with the same $(\corr,\ivar)$, which are Gaussians with the moments given by
\cref{eq:gc-forward-map,eq:gc-backward-map}. We prove by induction along the chain
that $\lmsg^{\mathrm p}_i=\lmsg^{\mathrm G}_i$ and $\rmsg^{\mathrm p}_i=\rmsg^{\mathrm G}_i$ for
every $i$.

\emph{Forward, base case.} $\lmsg_1=\mu=\normal{0}{1}$ in both chains; it is Gaussian,
so projecting it changes nothing.

\emph{Forward, inductive step.} Suppose $\lmsg^{\mathrm p}_{i-1}=\lmsg^{\mathrm G}_{i-1}=\normal{m}{v}$.
Multiplying by the Gaussian unary factor gives the tilted Gaussian $\normal{m'}{v'}$,
identical in the two chains. The transition integral in the non-Gaussian chain,
$\int\innov(a-\corr a')\normal{m'}{v'}(a')\dd a'$, is the density of $\corr U+V$ with
$U\sim\normal{m'}{v'}$ and $V\sim\innov$ independent; by \cref{lem:gc-conv} its mean is
$\corr m'$ and its variance $\corr^2v'+\ivar$, whatever $\innov$ is beyond its first two
moments. The projection therefore returns $\normal{\corr m'}{\corr^2v'+\ivar}$, which by
\cref{eq:gc-forward-map} is $\lmsg^{\mathrm G}_i$. The non-Gaussian shape of the integral
before projection (a mixture, a Laplace-smoothed Gaussian, whatever it may be) is
discarded by $\Pgauss$ and never influences a later step.

\emph{Backward, base case.} $\rmsg_\nsites\equiv1$ in both chains; it is not projected.
The first backward integral starts from the tilt of the flat message by
$\like(x_\nsites\mid\cdot)$, which is the Gaussian $\normal{x_\nsites/e^{-t}}{\Dt/e^{-2t}}$
in both chains.

\emph{Backward, inductive step.} Suppose $\rmsg^{\mathrm p}_{i+1}=\rmsg^{\mathrm G}_{i+1}$,
and let the tilt by $\like(x_{i+1}\mid\cdot)$ be $\normal{m'}{v'}$. The backward integral
in the non-Gaussian chain is $\int\innov(a''-\corr a)\normal{m'}{v'}(a'')\dd a''=h(\corr a)$
with $h(u)=\int\innov(a''-u)\normal{m'}{v'}(a'')\dd a''$, the density at $u$ of $U-V$ with
$U\sim\normal{m'}{v'}$, $V\sim\innov$; by \cref{lem:gc-conv} $h$ integrates to one and
has mean $m'$ and variance $v'+\ivar$. As a function of $a$, $h(\corr a)$ has mass
$1/|\corr|$, normalised mean $m'/\corr$ and normalised variance $(v'+\ivar)/\corr^2$
(substitute $u=\corr a$ in the moment integrals). Here $\corr\ne0$ is used: at
$\corr=0$ the function $h(0)$ is constant in $a$ and has no normalised moments. The
projection returns $\normal{m'/\corr}{(v'+\ivar)/\corr^2}$, which by
\cref{eq:gc-backward-map} is $\rmsg^{\mathrm G}_i$.

\emph{Beliefs.} At every site the projected belief is the product of three Gaussians,
$\lmsg^{\mathrm G}_i\like\rmsg^{\mathrm G}_i$, the belief of the Gaussian chain; its mean is
the Gaussian chain's posterior mean, which is the conditional mean of a jointly Gaussian
pair and hence, by \cref{eq:gc-lmmse-channel}, the LMMSE $e^{-t}\Sig_0\Sig_t^{-1}\vx$.
The recursion read $\innov$ only through $(0,\ivar)$, so the result is the same for
every such law. Nothing assumed stationarity: the induction carries site-dependent
$(m,v)$, and $\Sig_0$ is whatever covariance the chain has (with $a_1\sim\normal01$ and
$\ivar=1-\corr^2$ it is $\corr^{|i-j|}$, but the proposition does not need that).
\end{proof}

Two remarks on what is and is not checked by machine. The paper reports that the two
moment maps as integrals, the covariance identities $\Sig_{ax}=e^{-t}\Sig_0$,
$\Sig_{xx}=\Sig_t$, and the optimality of the normal-equation gain in arbitrary dimension
are verified in Lean~4 (\code{lean/README.md}), for a \emph{Gaussian} innovation; the
projection step and the moments-only clause of \cref{lem:gc-conv}, which are the
content of the proposition, are prose. That is an accurate description, and the reader
who wants to check the load-bearing step should check the two applications of
\cref{lem:gc-conv} above. Numerically, the proposition is checked in
\code{tests/test\_gaussian\_bp\_equivalence.py} by running a grid-mediated projection on
a Laplace chain and comparing with the information-form Gaussian recursion; the two
agree to the grid's accuracy at moderate $t$. The same file records an instability of
the grid-mediated projection at weakly informative $t$ (a nearly flat outgoing
message whose maximiser lies outside the domain is dragged to the boundary by the
truncation before projection), which is why the closed-form information recursion,
and not the grid projection, is what the paper uses for the LMMSE arm.

\section{Two consequences}

The paper draws two corollaries from the proposition ``as checks''. They are different
facts and each needs its own line.

\begin{corollary}[Every Gaussian-kernel sum-product is a linear map]
\label{cor:gc-linear}
Sum-product with a Gaussian transition kernel $\normal{\corr' a}{\ivar'}$ (any
$\corr',\ivar'$, not necessarily the truth) on the chain posterior returns posterior means
that are a fixed linear map $G(\corr',\ivar')\vx$. Consequently its mean-square error is at
least that of the LMMSE with the true covariance, and its score error is at least the
LMMSE floor $\E\lVert\score-\Sgauss\rVert^2$.
\end{corollary}
\begin{proof}
With a Gaussian kernel every message is Gaussian without any projection
(\cref{sec:gc-kalman}), the recursion is the Kalman smoother of the Gaussian chain with
parameters $(\corr',\ivar')$, and its posterior means are $\Sig'_{ax}\Sig'^{-1}_{xx}\vx$ for
the covariances of that chain: a linear map. By \cref{eq:gc-lmmse} the LMMSE with the
true second moments minimises the mean-square error over all linear maps, so it is a
floor for this one. For the score, Tweedie's identity gives
$\hat\score-\score=(e^{-t}/\Dt)(\hat\va-\pmean{\va})$ for any estimator $\hat\va$ of the
posterior mean, so the squared score error is $(e^{-2t}/\Dt^2)$ times the excess
mean-square error over the MMSE, and the ordering of estimators is the same.
\end{proof}

This is the sense in which the LMMSE is a \term{floor for every second-order
estimator}: a Gaussian-innovation kernel however fitted, a linear denoiser however
trained, and the Kalman smoother with any parameters all compute linear maps of $\vx$,
and none can beat $e^{-t}\Sig_0\Sig_t^{-1}\vx$. The paper uses the floor as a reference
line in every figure and table.

\begin{corollary}[The Gaussian quasi-likelihood sees only the covariance]
\label{cor:gc-quasi}
Fit the Gaussian chain with parameters $(\corr',\ivar')$ to noised data by maximising the
Gaussian marginal log-likelihood. That objective depends on the data only through the
sample covariance $\hat\Sig=\frac1\nseq\sum_\mu\vx^\mu\vx^{\mu\top}$ (per level), whose
expectation is $e^{-2t}\Sig_0+\Dt\Id$ whatever the innovation law; the population
maximiser is the true $(\corr,\ivar)$, so the fitted Gaussian kernel is consistent for
them and its sum-product lands on the LMMSE floor.
\end{corollary}
\begin{proof}
The Gaussian marginal log-likelihood of one sequence at level $t$ is
\[
  -\tfrac12\log\det(2\pi\Sig'_t)-\tfrac12\vx^{\!\top}\Sig'^{-1}_t\vx,\qquad
  \Sig'_t=e^{-2t}\Sig'_0+\Dt\Id;
\]
summed over sequences the data enter only through
$\sum_\mu\vx^{\mu\top}\Sig'^{-1}_t\vx^\mu=\nseq\tr(\Sig'^{-1}_t\hat\Sig)$. The expectation of
the per-sequence log-likelihood is
$-\tfrac12\log\det(2\pi\Sig'_t)-\tfrac12\tr(\Sig'^{-1}_t\Sig_t)$, which as a function of $\Sig'_t$ is maximised uniquely at $\Sig'_t=\Sig_t$ (the Gaussian
cross-entropy is minimised at the matching covariance); in the two-parameter family
$\Sig_t$ determines $(\corr,\ivar)$, since $\Sig_0=(\Sig_t-\Dt\Id)e^{2t}$ has entries
$\corr^{|i-j|}$ scaled by the variances determined by $\ivar$. Consistency of the
maximiser follows from the law of large numbers for $\hat\Sig$ and the continuity of the
objective.
\end{proof}

The distinction between the two corollaries is the distinction between the
proposition's two uses in the paper. \Cref{cor:gc-linear} says that the $\ncomp=1$ arm
\emph{cannot go below} the floor; \cref{cor:gc-quasi} says that it \emph{lands on} it,
which is what the capacity table and the regime sweep verify at every cell, and what
makes those cells an apparatus check of the whole pipeline (grid, sweeps, EM, scoring).

\section{The price of second-order closure}

Define the \term{closure gap} at level $t$ as $\E_{\Pt}\lVert\score-\Sgauss\rVert^2$, the
mean-square distance between the true score and that of the covariance twin, and the
\term{relative score error} of any estimator $\hat\score$ as
$\bigl(\sum_t\sum_{\mu,i}(\hat\score-\score)_i^2/\sum_t\sum_{\mu,i}\score_i^2\bigr)^{1/2}$,
pooled over the schedule (the paper's estimand, Appendix C). By
\cref{prop:gc-lmmse,cor:gc-linear} the closure gap is three things at once: the error of
Gaussian message passing on the non-Gaussian chain; the error of the best linear
denoiser; and the error of a Gaussian-innovation kernel fitted with infinite data. That
is why the paper calls it ``the price of second-order closure'' and ``what any estimator
that carries only second-order structure must leave on the table''. Its size is the
subject of the paper's order-by-order decay theorem (its Section 3): to leading order $\tfrac16\Frob{\kfour}^2e^{-8t}$,
the degree-four Hermite term the twin lacks, and the compendium's chapter on the
attenuation theorems derives it. Pooled over the twelve-level schedule and expressed
as a relative error it is the number the paper's Appendix D lists in the LMMSE column,
per innovation family.

The other reference point in that table is the \term{isotropic score} $-\vx$, the
score of $\normal{0}{\Id}$, which is the noised law of independent unit-variance sites.
It is the estimator that knows the marginal variance and nothing else. Through Tweedie,
$-\vx=-(\vx-e^{-t}\hat\va)/\Dt$ gives $\hat\va=(1-\Dt)\vx/e^{-t}=e^{-t}\vx$: the isotropic
denoiser shrinks the observation by $e^{-t}$, which is the LMMSE for $\Sig_0=\Id$, that
is, for a chain with $\corr=0$. Its relative error rises with $\corr$ because the
covariance itself is then most of what there is to learn, which is the paper's reading
of the figure in its Appendix F on the second-order floor across correlation. The
code for both references is \code{lmmse\_posterior\_mean},
\code{lmmse\_score} and \code{isotropic\_score} in \code{src/markov\_score/theory.py};
the isotropic posterior mean is written there as $(\vx-\Dt\vx)/e^{-t}$, the same thing.

\section{A worked example: two sites}

Take $\nsites=2$, $\corr=\tfrac12$, $\ivar=\tfrac34$ (so that $\Var a_2=\corr^2+\ivar=1$),
and a level with $e^{-t}=0.8$, $\Dt=0.36$. Observe $\vx=(1,0)$.

\emph{Closed form.} $\Sig_0=\begin{psmallmatrix}1&0.5\\0.5&1\end{psmallmatrix}$,
$\Sig_t=0.64\Sig_0+0.36\Id=\begin{psmallmatrix}1&0.32\\0.32&1\end{psmallmatrix}$ with
determinant $1-0.32^2=0.8976$, so
$\Sig_t^{-1}=\tfrac{1}{0.8976}\begin{psmallmatrix}1&-0.32\\-0.32&1\end{psmallmatrix}
=\begin{psmallmatrix}1.1141&-0.3565\\-0.3565&1.1141\end{psmallmatrix}$. Then
$\Sig_t^{-1}\vx=(1.1141,\,-0.3565)$, $\Sig_0\Sig_t^{-1}\vx=(1.1141-0.1783,\,0.5570-0.3565)=(0.9358,\,0.2005)$,
and the LMMSE is $0.8\times(0.9358,\,0.2005)=(0.7487,\,0.1604)$. The twin's score is
$-\Sig_t^{-1}\vx=(-1.1141,\,0.3565)$, and Tweedie's identity from the LMMSE gives the
same: $-(1-0.8\cdot0.7487)/0.36=-1.1141$ and $-(0-0.8\cdot0.1604)/0.36=0.3565$.

\emph{The two sweeps.} The unary factor has $\lambda_{\like}=0.64/0.36=1.7778$ and
$h_{\like}=0.8x_i/0.36$, that is $2.2222$ at site 1 and $0$ at site 2. Forward: the
prior at site 1 is $(\lambda,h)=(1,0)$; tilting gives $\lambda'=2.7778$, $m'=2.2222/2.7778=0.8$,
$v'=0.36$; the transition maps this to $m=0.5\cdot0.8=0.4$, $v=0.25\cdot0.36+0.75=0.84$:
$\lmsg_2=\normal{0.4}{0.84}$. Backward: the terminal message is flat; tilting it at
site 2 gives $m'=0$, $v'=1/1.7778=0.5625$; the inverse transition maps this to
$m=0/0.5=0$, $v=(0.5625+0.75)/0.25=5.25$: $\rmsg_1=\normal{0}{5.25}$. Belief at site 1:
precision $1+1.7778+1/5.25=2.9683$, information $0+2.2222+0=2.2222$, mean
$2.2222/2.9683=0.7487$. Belief at site 2: precision $1/0.84+1.7778=2.9683$, information
$0.4/0.84+0=0.4762$, mean $0.1604$. Both agree with the closed form, and the common
posterior variance $1/2.9683=0.3369$ is the diagonal of $\Sig_0-e^{-2t}\Sig_0\Sig_t^{-1}\Sig_0$.

\emph{A non-Gaussian innovation.} Replace $\innov$ by the two-point law
$\varepsilon=\pm\sqrt{0.75}$ with probability $\tfrac12$ each, which has the same mean and
variance. The exact forward message at site 2 is now the mixture
$\tfrac12\normal{0.4+0.866}{0.09}+\tfrac12\normal{0.4-0.866}{0.09}$, a bimodal shape whose
mean is $0.4$ and whose variance is $0.09+0.75=0.84$ (\cref{lem:gc-conv}). Its projection
is $\normal{0.4}{0.84}$, the Gaussian chain's message, and every later step is identical:
the projected recursion returns $(0.7487,\,0.1604)$ for this law too. The exact
posterior mean for the two-point law is different, and the difference is the closure
gap of this tiny example.

\where{The paper, Section 2 (``the covariance twin and the price of second-order
closure'', Proposition 1 and the two consequences), Appendix A (the proof), Section 3
(the size of the gap, the decay theorem), Section 5 and Appendices D and F (the LMMSE and
isotropic reference points, the $\ncomp=1$ arm landing on the floor in E02a, E02b, E13),
Appendix E (the Gaussian twin's windowed LMMSE in E15). Code:
\code{src/markov\_score/bp\_gaussian.py}, which holds \code{gaussian\_chain\_bp} and the
legacy grid-projected variant kept to document its instability;
\code{src/markov\_score/theory.py}, which holds \code{chain\_covariance},
\code{lmmse\_posterior\_mean}, \code{lmmse\_score} and \code{isotropic\_score};
\code{src/markov\_score/exact\_scores.py}; and \code{lean/README.md} for what is
machine-checked. Tests: \code{tests/test\_gaussian\_bp\_equivalence.py} and the
Gaussian-recovery test of \code{tests/test\_em\_bp.py}.}

\paragraph{References.}
R.~E.~Kalman (1960), ``A new approach to linear filtering and prediction problems'',
\emph{Journal of Basic Engineering} 82(1).
H.~E.~Rauch, F.~Tung and C.~T.~Striebel (1965), ``Maximum likelihood estimates of linear
dynamic systems'', \emph{AIAA Journal} 3(8).
T.~P.~Minka (2001), ``Expectation propagation for approximate Bayesian inference'',
\emph{Proceedings of UAI}.
S.~M.~Kay (1993), \emph{Fundamentals of Statistical Signal Processing: Estimation Theory},
Prentice Hall, Chapters 11 and 12.
C.~M.~Bishop (2006), \emph{Pattern Recognition and Machine Learning}, Springer,
Sections 2.3 and 13.3.
